\subsection{Datasets}

\textbf{Training Data} Our training corpus consists exclusively of the mathematics subset from the \texttt{guru-RL-92k} dataset, part of the Reasoning360 project\citep{cheng2025revisiting}. We directly leverage this pre-existing resource due to its rigorous curation process. The dataset's creators employed aggressive substring-based deduplication and a difficulty-based filtering method, which selected problems based on a baseline model's performance to ensure a high-quality and appropriately challenging training distribution. We further reorganized the dataset in the order of increasing difficulty (i.e., decreasing passrate), so as to enable the adoption of a curriculum learning paradigm for the training process.

\textbf{Evaluation Suite} Our evaluation is twofold. To derive scaling laws, we use a held-out set of 500 in-domain mathematics problems, sampled from the full dataset prior to training while preserving the original difficulty distribution. To assess generalization, we use a broader suite of benchmarks, covering mathematics, code generation, logical reasoning, and science, as detailed in Table~\ref{tab:eval_suite}.

\begin{table}[h]
\centering
\caption{Composition of the multi-domain evaluation suite.}
\label{tab:eval_suite}
\begin{tabular}{lrll}
\toprule
\textbf{Dataset} & \textbf{Samples} & \textbf{Data Source} & \textbf{Domain} \\
\midrule
Held-out Data & 500 & DeepScaler \& DAPO & Math \\
aime2024 & 30 & aime2024 & Math \\
amc2023 & 40 & aimeamc2023 & Math \\
codegen\_humaneval & 164 & codegen/humaneval & Code \\
gsm8k & 1319 & openai/gsm8k & Math \\
logic\_zebra\_puzzle & 200 & logic/zebra\_puzzle & Logical Reasoning \\
math & 500 & math/math & Math \\
stem\_supergpqa & 200 & stem/supergpqa & STEM \\
\midrule
\textbf{Total} & \textbf{2953} & & \\
\bottomrule
\end{tabular}
\end{table}